\documentclass[11pt]{article}
\usepackage[margin=1in]{geometry}
\usepackage{amsmath,amssymb,amsthm,mathtools}
\usepackage{graphicx}
\usepackage{hyperref}
\usepackage{xcolor}
\usepackage{enumitem}

\newtheorem{theorem}{Theorem}
\newtheorem{proposition}[theorem]{Proposition}
\newtheorem{lemma}[theorem]{Lemma}
\newtheorem{corollary}[theorem]{Corollary}
\theoremstyle{remark}
\newtheorem{remark}[theorem]{Remark}

\newcommand{\E}{\mathbb E}
\newcommand{\Pp}{\mathbb P}
\newcommand{\R}{\mathbb R}
\newcommand{\N}{\mathbb N}
\newcommand{\cL}{\mathcal L}
\newcommand{\Ran}{\operatorname{Ran}}
\newcommand{\HS}{\operatorname{HS}}

\title{No Logarithmic-Moment Criterion for Arbitrary\\
$C_0$-Semigroup L\'evy Convolutions}
\author{Automatic mathematics run \texttt{fa\_banach\_001}}
\date{17 August 2026}

\begin{document}
\setlength{\emergencystretch}{3em}
\maketitle

\begin{center}
\fbox{\parbox{0.92\textwidth}{
\textbf{Status.} Candidate counterexample, likely valid, for human review.
The logarithmic moment condition is neither necessary nor sufficient for
general strongly stable $C_0$-semigroups.  Necessity fails even for a
nilpotent exponentially stable contraction, and sufficiency fails even for
deterministic L\'evy drift and a bounded strongly stable contraction
generator.  Exact drift and Gaussian tests, an iid-block reduction, and the
sharp two-sided exponential regime are also proved.
}}
\end{center}

\section{Source question}

The source is Zbigniew J.~Jurek,
\emph{On Relations Between Urbanik and Mehler Semigroups},
arXiv:0811.2989 (2008) \cite{Jurek}.  On printed page 12 it asks whether,
for a $C_0$-semigroup $(T_t)_{t\geq0}$ on a separable Banach space and an
$E$-valued c\`adl\`ag L\'evy process $Y$, there is an analogue of the
bounded-generator criterion
\[
 \int_{(0,t]}e^{sV}\,dY(s)\ \text{converges as }t\to\infty
 \quad\Longleftrightarrow\quad
 \E\log(1+\|Y(1)\|)<\infty,
\]
or, more generally, how convergence depends on the infinitesimal generator
and on $Y(1)$.  The exact source passage is reproduced below.

\begin{center}
\includegraphics[width=0.97\textwidth]{figures/open_problem_crop.png}
\end{center}

The paper defines ``operator topology'' to mean the strong operator topology.
That convention matters: one of the examples below also shows that the
bounded-generator assertion immediately preceding the open question cannot
be correct with merely strong convergence.  A later review states the
classical iff with convergence in \emph{operator norm} \cite[Section
6.1]{Applebaum}.

\section{Answer and proof idea}

There are two independent obstructions.

\begin{itemize}[leftmargin=2em]
\item A $C_0$-semigroup may erase all old noise in finite time.  Then no tail
moment of the jumps is necessary.
\item Strong stability may be arbitrarily slow in different directions.  A
driver with the best possible moment behaviour can then accumulate in
infinitely many slowly decaying coordinates.
\end{itemize}

The first phenomenon is realized by the nilpotent translation semigroup on
$L^2(0,1)$ and heavy-tailed compound-Poisson noise.  The second is already
realized without randomness: on $\ell^2$, let
$T_te_n=e^{-t/n}e_n$ and drive the convolution by the deterministic L\'evy
process $Y(t)=t(1/n)_n$.

The correct structural replacement is a block series.  Writing the
convolution over successive unit intervals produces iid variables $\xi_n$
and the exact identity
\[
 \int_{(0,N]}T_s\,dY(s)=\sum_{n=0}^{N-1}T_n\xi_{n+1}.
\]
Thus the generator contributes its discrete orbit $(T_n)$, while the law of
$Y(1)$ contributes the one-block infinitely divisible law.  This reduction
gives closed criteria for deterministic drift and Gaussian noise.  It also
shows exactly when a logarithmic test returns: an exponential upper estimate
gives sufficiency, and an exponential lower estimate gives necessity.

\section{The logarithmic condition fails in both directions}

\begin{theorem}[Two-sided failure]\label{thm:two-sided}
There are separable Hilbert spaces, strongly stable contraction
$C_0$-semigroups, and c\`adl\`ag L\'evy processes for which each of the
following holds.
\begin{enumerate}[label=\textup{(\roman*)}]
\item The improper convolution converges almost surely although
$\E\log(1+\|Y(1)\|)=\infty$.  The semigroup may be chosen nilpotent and hence
exponentially stable.
\item $Y(1)$ is deterministic (hence has every finite moment), but the
convolution norm tends to infinity.  The generator may be chosen bounded and
the semigroup strongly stable and contractive.
\end{enumerate}
Consequently no logarithmic-moment iff criterion extends to arbitrary
strongly stable $C_0$-semigroups.  Under the source paper's explicit strong
operator convention, part \textup{(ii)} also corrects its preceding
bounded-generator formulation.
\end{theorem}

\begin{proof}
For \textup{(i)}, take $E=L^2(0,1)$ and extend every $f\in E$ by zero outside
$(0,1)$.  Set
\[
 (T_tf)(r)=f(r+t),\qquad 0<r<1.
\]
Translations on the zero extension show strong continuity; $(T_t)$ is a
contraction semigroup and $T_t=0$ for $t\geq1$.

Fix a unit vector $v\in E$.  Let $Y$ be a rate-one compound-Poisson process
whose jumps are $R_kv$, where the positive iid amplitudes satisfy
\[
 \Pp\{R=e^m\}=\frac{6}{\pi^2m^2},\qquad m\geq1.
\]
If $N_1$ is the number of jumps by time one, then on $\{N_1=1\}$ one has
$Y(1)=Rv$.  Therefore
\begin{align*}
 \E\log(1+\|Y(1)\|)
 &\geq \Pp(N_1=1)\E\log(1+R)\\
 &\geq e^{-1}\frac{6}{\pi^2}\sum_{m\geq1}\frac{m}{m^2}
 =\infty.
\end{align*}
On the other hand, if $S_k$ are the jump times, then
\[
 \int_{(0,t]}T_s\,dY(s)
 =\sum_{S_k\leq \min\{t,1\}}T_{S_k}(R_kv).
\]
This is a finite sum and is constant for $t\geq1$, so it converges almost
surely.

For \textup{(ii)}, take $E=\ell^2$ with standard basis $(e_n)$.  Define
\[
 Ve_n=-\frac1n e_n,\qquad T_t=e^{tV},\qquad
 T_te_n=e^{-t/n}e_n.
\]
The operator $V$ is bounded and $(T_t)$ is a contraction semigroup.  For each
$x=(x_n)\in\ell^2$, dominated convergence gives
\[
 \|T_tx\|^2=\sum_{n\geq1}e^{-2t/n}|x_n|^2\longrightarrow0,
\]
so $T_t\to0$ strongly.  It does not converge in operator norm, since
$\|T_t\|=\sup_n e^{-t/n}=1$.

Let $y=(1/n)_{n\geq1}\in\ell^2$ and let $Y(t)=ty$.  This is a deterministic
L\'evy process and $\E\log(1+\|Y(1)\|)<\infty$.  Its convolution is the
Bochner integral
\[
 X_t=\int_0^tT_sy\,ds
     =\bigl(1-e^{-t/n}\bigr)_{n\geq1}.
\]
For every fixed $t$, this vector belongs to $\ell^2$ because
$1-e^{-t/n}\leq t/n$.  But monotone convergence yields
\[
 \|X_t\|^2=\sum_{n\geq1}(1-e^{-t/n})^2\longrightarrow
 \sum_{n\geq1}1=\infty.
\]
Thus the deterministic laws $\delta_{X_t}$ are not tight, and there is no
convergence in distribution, in probability, or almost surely.
\end{proof}

\begin{remark}[Why exponential stability only repairs sufficiency]
The nilpotent example satisfies $\|T_t\|=0$ for $t\geq1$, so even exponential
stability cannot make logarithmic integrability necessary.  The diagonal
example has only strong stability; this is exactly why the usual exponential
upper-bound proof of sufficiency is unavailable.
\end{remark}

\section{Exact generator criterion for deterministic drift}

The deterministic witness is not an isolated trick; it is the complete drift
case.

\begin{theorem}[Range criterion]\label{thm:range}
Let $(T_t)$ be a strongly stable $C_0$-semigroup on a Banach space $E$, with
generator $A$, and let $Y(t)=ty$.  Then
\[
 \int_0^tT_s\,dY(s)=\int_0^tT_sy\,ds
\]
converges in norm as $t\to\infty$ if and only if
\[
 y\in\Ran(-A).
\]
If $y=-Ax$, the limit is $x$, and this $x\in D(A)$ is unique.
\end{theorem}

\begin{proof}
If $y=-Ax$ with $x\in D(A)$, the fundamental semigroup identity gives
\[
 \int_0^tT_sy\,ds=x-T_tx\longrightarrow x.
\]
Conversely, suppose
$x_t:=\int_0^tT_sy\,ds\to x$ in norm.  The integrated-semigroup identity
says
\[
 x_t\in D(A),\qquad Ax_t=T_ty-y.
\]
Strong stability makes the right side converge to $-y$.  Since $A$ is
closed, $x\in D(A)$ and $Ax=-y$.  Finally, if $Az=0$, then $T_tz=z$ for all
$t$, while strong stability gives $T_tz\to0$; hence $z=0$.
\end{proof}

For the diagonal generator in Theorem~\ref{thm:two-sided},
$(-A)x=(x_n/n)_n$.  The chosen $y=(1/n)_n$ would require
$x=(1,1,\ldots)\notin\ell^2$, so $y\notin\Ran(-A)$, exactly as the theorem
predicts.

\section{Exact iid-block reduction}

Assume throughout this section that the finite-horizon stochastic integrals
exist in $E$.  Put
\[
 X_t:=\int_{(0,t]}T_s\,dY(s).
\]

\begin{theorem}[Block-series criterion]\label{thm:block}
For $n\geq0$, define
\[
 \xi_{n+1}:=\int_{(n,n+1]}T_{s-n}\,dY(s).
\]
Then $(\xi_n)_{n\geq1}$ is iid, with common law
\[
 \cL(\xi)=\cL\left(\int_{(0,1]}T_s\,dY(s)\right),
\]
and, pathwise for every integer $N\geq1$,
\[
 X_N=\sum_{n=0}^{N-1}T_n\xi_{n+1}.                 \tag{4.1}
\]
If $\Psi$ is the characteristic exponent of $Y(1)$, then the one-block law
is determined by
\[
 \log\E e^{i\langle x^*,\xi\rangle}
   =\int_0^1\Psi(T_s^*x^*)\,ds,
 \qquad x^*\in E^*.                                \tag{4.2}
\]

Suppose in addition that $T_n\to0$ strongly and that the one-block
convolution process has uniformly tight marginals on $[0,1]$.  Then
convergence of $X_t$ in probability as $t\to\infty$ is equivalent to
convergence in probability of the series in \textup{(4.1)}.
\end{theorem}

\begin{proof}
The processes
$u\mapsto Y(n+u)-Y(n)$ on disjoint unit intervals are independent and have
the same law, which proves the iid assertion.  Bounded deterministic
operators commute with the stochastic integral, and the semigroup law gives
\[
 T_n\xi_{n+1}
 =\int_{(n,n+1]}T_nT_{s-n}\,dY(s)
 =\int_{(n,n+1]}T_s\,dY(s).
\]
Summing proves (4.1).  Formula (4.2) follows first for step-function
integrands from independent increments and then for the defining
finite-horizon integral by approximation.

For the last assertion write $t=N+r$, $0\leq r<1$.  Then
\[
 X_t-X_N=T_N\xi_{N+1}(r),
 \quad
 \xi_{N+1}(r):=\int_{(N,N+r]}T_{s-N}\,dY(s).
\]
The family $\{\xi_{N+1}(r_N):N\geq1\}$ is tight for every deterministic
choice $r_N\in[0,1]$.  Strong convergence $T_N\to0$ is uniform on compact
subsets of $E$: this follows from the uniform boundedness principle and a
finite-net argument.  Hence $T_N\xi_{N+1}(r_N)\to0$ in probability.  Apply
this to arbitrary real-time sequences and use the integer subsequence for
the converse.
\end{proof}

This is an exact general answer at the structural level.  A $C_0$-generator
determines $(T_t)$ uniquely, while $Y(1)$ determines the entire L\'evy
process law; equation (4.2) displays their interaction rather than reducing
it to one radial moment.

\section{The sharp exponential regime for a logarithmic test}

\begin{theorem}[Two-sided exponential criterion]\label{thm:log}
Let $(\xi_n)$ be iid $E$-valued random variables.  Suppose that for constants
$c,C,\omega,\Omega>0$,
\[
 c e^{-\Omega n}\|x\|\leq\|T_nx\|
 \leq C e^{-\omega n}\|x\|,
 \qquad n\geq0,\ x\in E.                            \tag{5.1}
\]
Then
\[
 \sum_{n\geq0}T_n\xi_{n+1}\ \text{converges almost surely}
 \quad\Longleftrightarrow\quad
 \E\log(1+\|\xi_1\|)<\infty.                      \tag{5.2}
\]
Only the upper estimate in (5.1) is needed for the forward implication from
the logarithmic moment to convergence; only the lower estimate is needed for
the reverse implication.
\end{theorem}

\begin{proof}
Let $Z=\log(1+\|\xi_1\|)$.  If $\E Z<\infty$, then for every
$\varepsilon>0$,
\[
 \sum_{n\geq1}\Pp\{Z>\varepsilon n\}<\infty
\]
by comparison with the tail integral for $\E Z$.  Borel--Cantelli gives
$\|\xi_n\|\leq e^{\varepsilon n}$ eventually.  Choose
$0<\varepsilon<\omega$; the upper bound in (5.1) then makes the series
absolutely convergent almost surely.

Conversely, almost-sure convergence forces
$\|T_n\xi_{n+1}\|\to0$.  The lower bound implies, eventually,
\[
 \|\xi_{n+1}\|\leq c^{-1}e^{\Omega n}.
\]
If $\E Z=\infty$, comparison with the tail integral instead gives
\[
 \sum_{n\geq1}
 \Pp\{\|\xi_{n+1}\|>c^{-1}e^{\Omega n}\}=\infty.
\]
These events are independent, so the second Borel--Cantelli lemma says they
occur infinitely often, a contradiction.
\end{proof}

The two estimates isolate the counterexamples precisely.  The nilpotent
shift has no exponential lower bound; the diagonal semigroup has no uniform
decay upper bound.  If $A$ is bounded, then $T_t=e^{tA}$ extends to a group and
\[
 \|T_nx\|\geq\|e^{-nA}\|^{-1}\|x\|
 \geq e^{-n\|A\|}\|x\|.                             \tag{5.3}
\]
Thus a bounded generator automatically supplies the missing lower bound.  If
in addition $\|T_t\|\to0$, the semigroup law upgrades this to an exponential
upper bound.

For the block variable in Theorem~\ref{thm:block}, boundedness of $A$ also
makes $T_s$ and $T_s^{-1}$ uniformly bounded for $0\leq s\leq1$.  Its L\'evy
measure is
\[
 \nu_\xi(B)=\int_0^1\nu\bigl(T_s^{-1}B\bigr)\,ds.
\]
The standard logarithmic-moment criterion for infinitely divisible laws
therefore shows
\[
 \E\log(1+\|\xi\|)<\infty
 \quad\Longleftrightarrow\quad
 \E\log(1+\|Y(1)\|)<\infty.                         \tag{5.4}
\]
Equations (5.1)--(5.4) recover the classical bounded-generator criterion at
integer times; the standard one-block maximal estimate gives its all-time
form.  In particular, they show why \emph{operator-norm} stability and the
bounded-generator lower estimate are the two decisive features missing from
the general $C_0$ question.

\section{Gaussian noise: an exact covariance test}

For completeness, the other principal noise component has its own exact
criterion.  This is a standard Hilbert-space fact, included to make clear why
a radial moment cannot encode the answer.

\begin{proposition}[Gaussian energy criterion]\label{prop:gauss}
Let $H$ be a separable Hilbert space and let $Y=W_Q$ be a $Q$-Wiener process,
where $Q$ is positive and trace class.  Then
\[
 \int_0^tT_s\,dW_Q(s)
\]
converges in $L^2(\Omega;H)$ as $t\to\infty$ if and only if
\[
 \int_0^\infty\|T_sQ^{1/2}\|_{\HS}^2\,ds<\infty.   \tag{6.1}
\]
\end{proposition}

\begin{proof}
For $0\leq u<t$, the It\^o isometry gives
\[
 \E\left\|\int_u^tT_s\,dW_Q(s)\right\|^2
 =\int_u^t\|T_sQ^{1/2}\|_{\HS}^2\,ds.
\]
Thus the convolutions are Cauchy in $L^2(\Omega;H)$ exactly when (6.1)
holds.
\end{proof}

In the diagonal semigroup above, take
$W_Q(t)=\sum_{n\geq1}n^{-1}\beta_n(t)e_n$.  Then $W_Q(1)$ is Gaussian and
has every finite moment, but
\[
 \int_0^\infty\|T_sQ^{1/2}\|_{\HS}^2ds
 =\frac12\sum_{n\geq1}\frac1n=\infty.
\]
More strongly, the convolution coordinates have variances
\[
 \sigma_n^2(t)=\frac{1-e^{-2t/n}}{2n}.
\]
Writing $M_t=\sum_n\sigma_n^2(t)$, one has $M_t\to\infty$ and
\[
 \operatorname{Var}(\|X_t\|^2)
 =2\sum_n\sigma_n^4(t)
 \leq\frac12\sum_n\frac1{n^2}.
\]
Chebyshev's inequality yields $\|X_t\|\to\infty$ in probability.  Hence even
Gaussian noise with all moments can fail to have a limiting convolution under
mere strong stability.

\section{Literature boundary, limitations, and review recommendation}

Applebaum's review \cite[Theorem 6.3 and Section 6.1]{Applebaum} records the
known Hilbert-space L\'evy-triplet characterization, attributed to
Chojnowska--Michalik: Gaussian trace energy, integrated L\'evy-measure mass,
and a compensated-drift limit.  It also records logarithmic integrability as
sufficient under exponential stability, and the iff result for a contraction
semigroup with bounded generator converging in operator norm.  Those known
results are consistent with every statement above.

The run's registry and attempts indexes were searched for arXiv:0811.2989 and
for the central phrases.  A bounded arXiv/web search through 17 August 2026
used the exact title, the exact question, ``$C_0$ semigroup stochastic
convolution logarithmic moment'', ``bounded generator'', and
``Chojnowska--Michalik''.  It found the source and Applebaum's review, but no
later paper explicitly claiming to answer Jurek's general Banach-space
question.  Novelty is therefore plausible, not certified.

The packet does \emph{not} claim a new arbitrary-Banach L\'evy-triplet
classification, nor does it claim the Hilbert triplet theorem as new.  Its
complete claim is the negative answer to a universal logarithmic analogue,
the correction of the source's strong-topology formulation, and the exact
replacement statements proved here.

Human review should focus on three short points: the topology convention in
the source, the block-law logarithmic equivalence (5.4), and the stated
finite-horizon tightness condition in Theorem~\ref{thm:block}.  The two
counterexamples and Theorem~\ref{thm:range} are elementary and independent of
those auxiliary points.

\begin{thebibliography}{9}

\bibitem{Jurek}
Z.~J. Jurek,
\emph{On Relations Between Urbanik and Mehler Semigroups},
arXiv:0811.2989, 2008.
\url{https://arxiv.org/abs/0811.2989}.

\bibitem{Applebaum}
D.~Applebaum,
\emph{Infinite Dimensional Ornstein--Uhlenbeck Processes Driven by L\'evy
Processes},
arXiv:1411.2752, 2014.
\url{https://arxiv.org/abs/1411.2752}.

\bibitem{CM}
A.~Chojnowska--Michalik,
\emph{On processes of Ornstein--Uhlenbeck type in Hilbert space},
Stochastics 21 (1987), 251--286.

\end{thebibliography}

\end{document}
